\chapter{减速器系统核心理论}
\label{ch:gearbox}

\section{减速器的本质——力矩放大器与误差放大器}

减速器在机器人关节中的核心作用只有一个：\textbf{将电机的高速小力矩转换为低速大力矩}。

\begin{equation}
T_{\text{out}} = T_{\text{motor}} \times i \times \eta
\qquad
\omega_{\text{out}} = \frac{\omega_{\text{motor}}}{i}
\end{equation}

其中$i$是减速比，$\eta$是传动效率（$0 < \eta < 1$）。

\begin{keyinsight}{减速器也是误差放大器}
这是本书最重要的认知之一：\\
减速器在放大\textbf{力矩}的同时，也在放大\textbf{所有非线性误差}。\\
放大倍数 = 减速比 $i$ 本身。
\end{keyinsight}

\section{四种主流减速器结构}

\subsection{行星减速器}

\begin{itemize}
  \item \textbf{结构}：太阳轮+行星轮+齿圈
  \item \textbf{减速比范围}：3:1 -- 100:1（单级3--10，多级可达100）
  \item \textbf{效率}：90--97\%（单级），多级略低
  \item \textbf{特点}：结构紧凑、刚度高、背隙小、成本适中
  \item \textbf{应用}：四足机器人、协作机器人、工业机器人小臂
  \item \textbf{控制影响}：背隙（backlash）在换向时产生力矩死区
\end{itemize}

\subsection{谐波减速器}

\begin{itemize}
  \item \textbf{结构}：波发生器+柔轮+刚轮
  \item \textbf{减速比范围}：30:1 -- 160:1
  \item \textbf{效率}：70--85\%
  \item \textbf{特点}：单级减速比大、零背隙、重量轻
  \item \textbf{应用}：人形机器人协作臂、精密定位
  \item \textbf{控制影响}：柔轮扭转刚度低，引入低频谐振；摩擦非线性强
\end{itemize}

\begin{cautionbox}{谐波减速器的"扭转柔性陷阱"}
谐波减速器的柔轮在大力矩下会产生\textbf{显著扭转形变}（可达1--3角分/Nm）。\\
这个柔性在力控中表现为：\\
你发了力矩指令，实际输出要\textbf{等柔轮扭到对应角度}力矩才建立。\\
这就是\textbf{力控相位滞后}的来源之一。
\end{cautionbox}

\subsection{直驱（Direct Drive）}

\begin{itemize}
  \item \textbf{结构}：电机直接连接负载，无减速器
  \item \textbf{减速比}：1:1
  \item \textbf{效率}：接近100\%
  \item \textbf{特点}：零背隙、零摩擦（减速器部分）、响应极快
  \item \textbf{应用}：高精度力矩控制、直接力控
  \item \textbf{控制影响}：优点是力控带宽最高；缺点是需要巨大扭矩电机
\end{itemize}

\subsection{半直驱（Quasi-Direct Drive, QDD）}

\begin{itemize}
  \item \textbf{结构}：低减速比（通常6:1 -- 15:1）+ 大扭矩电机
  \item \textbf{效率}：85--95\%
  \item \textbf{特点}：介于直驱和传统减速器之间，兼顾力矩密度和柔顺性
  \item \textbf{应用}：宇树H1、MIT Cheetah、四足机器人
  \item \textbf{控制影响}：这是\textbf{最合适力控的折中方案}，详见第7章
\end{itemize}

\begin{table}[H]
\centering
\begin{tabular}{lcccc}
\toprule
\textbf{类型} & \textbf{减速比} & \textbf{效率} & \textbf{背隙} & \textbf{力控适合度} \\
\midrule
行星 & 3--100 & 90--97\% & 有 & ★★★☆☆ \\
谐波 & 30--160 & 70--85\% & 无（但有柔性） & ★★☆☆☆ \\
直驱 & 1 & $\approx$100\% & 无 & ★★★★★ \\
半直驱 & 6--15 & 85--95\% & 极小 & ★★★★☆ \\
\bottomrule
\end{tabular}
\caption{四种减速器类型对比}
\label{tab:gearbox_types}
\end{table}

\section{减速比对控制的核心影响——误差放大公式}

这是本书\textbf{最重要的公式}，请务必理解：

\subsection{电机端误差的传递}

假设电机端有一个误差源$\varepsilon_{\text{motor}}$（例如齿槽转矩、电流噪声、$K_t$波动），经过减速器后，输出端的误差为：

\begin{equation}
\varepsilon_{\text{out}} = \varepsilon_{\text{motor}} \times i
\end{equation}

\textbf{减速比$i$的物理意义}：它同时是力矩放大倍数和误差放大倍数。

\subsection{具体量化例子}

以$i=121$（你躯干关节的减速比）为例：

\begin{table}[H]
\centering
\begin{tabular}{lcc}
\toprule
\textbf{误差源} & \textbf{电机端幅值} & \textbf{关节输出端幅值（$i=121$）} \\
\midrule
齿槽转矩 & 0.01 Nm & \textbf{1.21 Nm} \\
电流噪声（1\%满量程） & 0.005 Nm & \textbf{0.605 Nm} \\
$K_t$波动（3\%） & 力矩偏移3\% & \textbf{力矩偏移3\%（相对放大）} \\
轴承摩擦波动 & 0.005 Nm & \textbf{0.605 Nm} \\
\bottomrule
\end{tabular}
\caption{$i=121$时误差放大效应}
\label{tab:error_amplification}
\end{table}

\begin{cautionbox}{为什么i=121无法做电流力矩观测}
因为电机端的微小误差（比如0.01 Nm的齿槽转矩）\\
在关节端被放大到\textbf{1.21 Nm}。\\
而你的关节最大力矩可能只有几十Nm。\\
力矩观测的信噪比\textbf{完全被减速比的误差放大效应摧毁}。
\end{cautionbox}

\section{三大致命减速比区间}

根据减速比的大小，机器人关节被分为\textbf{三个控制特性完全不同的区间}：

\subsection{区间一：低减速比（$i \leq 25$）——可电流力矩观测}

\begin{itemize}
  \item 误差放大倍数小，电机端误差在可接受范围
  \item 电流$\rightarrow$力矩的线性度好
  \item 可以实现\textbf{无力矩传感器的力控}
  \item 代表：宇树H1（$i\approx 10$）、MIT Cheetah（$i\approx 8$）
\end{itemize}

\subsection{区间二：中减速比（$25 < i < 80$）——勉强可用}

\begin{itemize}
  \item 误差已经明显，但可通过摩擦力补偿+前馈表部分抵消
  \item 力控精度受限，但仍然可以做一些基本的柔顺控制
  \item 代表：部分协作机器人关节
\end{itemize}

\subsection{区间三：高减速比（$i \geq 80$）——电流力矩观测彻底失效}

\begin{itemize}
  \item 误差放大倍数过大，电流观测的信号被噪声淹没
  \item \textbf{必须使用输出端力矩传感器}才能做力控
  \item 代表：你当前的躯干关节（$i=121$）、特斯拉Optimus（$i\approx 100$）
\end{itemize}

\begin{figure}[H]
\centering
\begin{tcolorbox}[colback=lightgray]
\centering
\textbf{减速比与控制可行性图谱}\\
\vspace{5pt}
$i \leq 25$：\hfill \textcolor{successgreen}{■ 电流力控可行} \hfill 宇树H1\\
$25 < i < 80$：\hfill \textcolor{infoyellow}{■ 勉强可用} \hfill 部分协作臂\\
$i \geq 80$：\hfill \textcolor{warnred}{■ 电流力控失效} \hfill 特斯拉/躯干关节\\
\end{tcolorbox}
\caption{减速比区间与电流力控可行性}
\label{fig:ratio_zones}
\end{figure}

\section{减速器的其他非线性——摩擦、背隙、柔性}

\subsection{摩擦}

减速器引入的摩擦力矩通常远大于电机本身的摩擦力矩：

\begin{equation}
T_{\text{friction, gearbox}} \approx 5\text{--}20\% T_{\text{rated}}
\end{equation}

摩擦特性包括：
\begin{itemize}
  \item \textbf{Stribeck效应}：低速时摩擦随速度增加而减小——这是低速力控震荡的物理原因
  \item \textbf{静摩擦+动摩擦差异}：启动时需要的力矩远大于维持运动需要的力矩
  \item \textbf{温度依赖性}：冷启动摩擦比热机大2--3倍
\end{itemize}

\subsection{背隙（Backlash）}

\begin{itemize}
  \item 主要存在于行星减速器中，谐波/直驱/半直驱基本无背隙
  \item 背隙导致\textbf{力矩死区}：在换向时，电机转了一个角度但输出端不动
  \item 背隙会\textbf{破坏位置控制的精度}和\textbf{力矩控制的连续性}
\end{itemize}

\subsection{扭转柔性}

\begin{itemize}
  \item 谐波减速器的柔轮扭转刚度$K_{\text{harmonic}} \approx 10^3\text{--}10^4$ Nm/rad
  \item 在大力矩下，扭转形变可达0.5--2度
  \item 这引入了\textbf{额外的谐振模态}，降低力控带宽
\end{itemize}

\section{减速比选择的工程权衡}

选择减速比不是越大越好，需要在以下因素之间权衡：

\begin{table}[H]
\centering
\begin{tabular}{p{5cm}|p{5cm}}
\toprule
\textbf{增大减速比的好处} & \textbf{增大减速比的代价} \\
\midrule
\begin{minipage}{4.8cm}
\vspace{2pt}
$\bullet$ 电机力矩要求降低（更小/更轻的电机）\\
$\bullet$ 更高的位置分辨率\\
$\bullet$ 更好的自锁能力\\
\vspace{2pt}
\end{minipage}
&
\begin{minipage}{4.8cm}
\vspace{2pt}
$\bullet$ 误差放大\\
$\bullet$ 力控灵敏度降低\\
$\bullet$ 摩擦和效率损失增加\\
$\bullet$ 输出端惯量反射降低（对外力不敏感）\\
\vspace{2pt}
\end{minipage} \\
\bottomrule
\end{tabular}
\caption{减速比选择的权衡}
\label{tab:ratio_tradeoff}
\end{table}

\begin{engineeringbox}{给硬件团队提需求的标准}
当你需要为某个关节选择减速比时，先问自己一个问题：\\
\textbf{这个关节需要做力控吗？}\\
\begin{itemize}
  \item 如果不能力控 → 选高减速比（便宜、定位准）
  \item 如果需要力控 → 减速比必须 $\leq 25$
  \item 如果必须高减速比又需要力控 → \textbf{必须加输出端力矩传感器}
\end{itemize}
\end{engineeringbox}

\section*{本章小结}
\addcontentsline{toc}{section}{本章小结}

\begin{summarybox}{}
\begin{enumerate}
  \item 减速器是\textbf{力矩放大器也是误差放大器}，放大倍数=$i$
  \item 三大减速比区间：$i \leq 25$（电流力控可行）、$25 < i < 80$（勉强可用）、$i \geq 80$（彻底失效）
  \item 摩擦力矩+背隙+扭转柔性是减速器三大非线性源
  \item \textbf{$i=121$的关节，电机端0.01 Nm的误差在关节端变成1.21 Nm}
  \item 这就是为什么你躯干关节（$i=121$）\textbf{做不了电流力矩观测}——物理限制，不是算法问题
\end{enumerate}
\end{summarybox}

\newpage
